 
% ============================================================
% Inflation Broadening in the Tails — DRAFT SKELETON
% Annotated outline: section structure + key points to discuss.
% Bullets describe WHAT to write where; established results are
% embedded so this doubles as a writing roadmap.
% Headline spec: 6 BLS CPI subcomponents (no oil); R2QConnectedness
% (nearPD); Schaefer-Strimmer shrinkage; nlag=2 main / nlag=1 robustness;
% tau in {0.1,0.3,0.5,0.7,0.9}; full sample Feb 1967 - May 2026.
% Compiler: pdfLaTeX + biber.
% ============================================================
\documentclass{article}

\usepackage[T1]{fontenc}
\usepackage[utf8]{inputenc}
\usepackage{mathpazo}
\usepackage[a4paper, total={6in, 8in}]{geometry}
\usepackage{amsmath,amssymb}
\usepackage{graphicx}
\usepackage{booktabs}
\usepackage{array,tabularx,multirow}
\usepackage{threeparttable,threeparttablex,longtable}
\usepackage{dcolumn}
\usepackage{makecell}
\usepackage{rotating}
\usepackage{pdflscape}
\usepackage{subcaption}
\usepackage{ragged2e}
\usepackage{xcolor}
\usepackage{float}
\usepackage{caption}
\usepackage{footnote}
\usepackage{authblk}
\usepackage[tableposition=top]{caption}
\usepackage{setspace}
\usepackage{bm}

\usepackage[style=apa,backend=biber,uniquename=false,natbib=true]{biblatex}
\usepackage[autostyle=true]{csquotes}
\addbibresource{paper/references.bib}

\makesavenoteenv{tabular}
\makesavenoteenv{table}

\usepackage[hidelinks]{hyperref}

\captionsetup[figure]{labelfont=bf,name=Fig.,labelsep=space}
\captionsetup[table]{labelfont=bf,name=Table,labelsep=space}

\onehalfspacing
% NOTE: this is a STRUCTURE/OUTLINE draft. Bullets = what to write.
\newcommand{\todo}[1]{\textcolor{red}{[\textbf{TODO:} #1]}}

\title{When prices move together: connectedness among CPI components, wages, and expectations across the inflation distribution and its consequences}

\author[a]{Michael Ryan}
\affil[a]{School of Accounting, Finance and Economics, University of Waikato, Hamilton, New Zealand.
Email: michael.ryan@waikato.ac.nz}
\date{Preliminary draft --- \today}

\begin{document}
\maketitle


\begin{abstract}
 
 
 
\noindent\textbf{Keywords:} inflation; CPI components; quantile connectedness; broadening; tail dependence. \\
\noindent\textbf{JEL:}  
\end{abstract}

% ============================================================
\section{Introduction}\label{sec:intro}
% ============================================================

Inflation targeting has made inflation one of the most closely monitored
macroeconomic indicators. When inflation rises sharply, attention shifts from
headline inflation alone to the breadth, persistence and propagation of price
increases. In 2021 at Jackson Hole, \citet{Powell2021} played down the rise in
inflation as narrow and likely temporary, describing it as ``largely the product
of a relatively narrow group of goods and services that have been directly
affected by the pandemic.'' By 2022, that assessment had changed:
\citet{Powell2022JH} observed that ``high inflation has continued to spread
through the economy'' and, drawing on the lessons of the 1970s, warned that the
anticipation of high inflation can become ``entrenched in the economic
decisionmaking of households and businesses.'' 

Powell's remarks suggest two questions become central for monetary policy when inflation rises. First, are price
increases concentrated in a small number of categories, or have they become
broad-based across the household expenditure basket? A relative price shock affecting a
limited set of goods can often be looked through; a generalised increase across
many categories is harder to dismiss \citep{bryan1994measuring}. Second, has
inflation become self-reinforcing through second-round effects in wages and
inflation expectations? If higher prices raise wage demands and inflation
expectations, and those wages and expectations then feed back into prices,
inflation may become more persistent and costly to reduce.

% ============================================================
% RELATED LITERATURE  (to be written)
%   Situate the paper in the inflation-dynamics literature before the
%   contributions paragraph. Suggested cites by theme are listed below;
%   delete this block once the review is written.
% -- Breadth / core / underlying inflation: bryan1994measuring;
%    Bryan-Cecchetti-Wiggins (1997); Stock-Watson (2007, 2016);
%    Reis-Watson (2010).
% -- The 2021-22 surge: Ball-Leigh-Mishra (2022, BPEA); Shapiro (2022, FRBSF);
%    Blanchard-Bernanke (2023); di Giovanni et al. (2023).
% -- The Great Inflation: DeLong (1997); Bordo-Orphanides (2013);
%    Primiceri (2006); Sargent-Williams-Zha (2006); Orphanides (2001).
% -- Expectations (adaptive vs anchored): Coibion-Gorodnichenko-Kamdar (2018, JEL);
%    Coibion-Gorodnichenko (2015); Ball-Mazumder (2011).
% -- Wage-price dynamics: Alvarez et al. (2022, IMF); Lorenzoni-Werning (2023).
% -- Disaggregated inflation / transmission: Boivin-Giannoni-Mihov (2009);
%    Pasten-Schoenle-Weber (2020); Ciccarelli-Mojon (2010).
% -- Connectedness methods: dieboldyilmaz2012; antonakakis2020; balli2023;
%    ando2022; stenfors2026decomposing; shahzadetal.
% ============================================================
\todo{[lit review here]}

This paper examines whether high inflation changes the connectedness---the strength
of statistical association---among prices, wages and inflation expectations in the
United States, and whether changes in price, wage and inflations expectation connectedness has consequences for the conduct of
monetary policy. Connectedness measures, developed to quantify directional spillovers
in financial systems \autocite{dieboldyilmaz2012}, are well suited to this question
because they provide a direct measure of how movements in one series are associated
with movements in others. The closest prior contribution is
\textcite{stenfors2026decomposing}, who study connectedness among disaggregated CPI
series for the United States, the United Kingdom and Japan since the late 1980s using
the time-varying-parameter VAR (TVP-VAR) framework of \textcite{antonakakis2020}. We
make three contributions relative to that work. The first is \emph{span}: we extend
the sample back to the Great Inflation, so that the high inflation of the 1970s and the
COVID-era  high inflation of 2021--22 can be compared within a common empirical framework. The
second is \emph{scope}: alongside the six disaggregated CPI categories we include
wages and inflation expectations, which lets us examine wage--price interactions and the
inflation--expectations interactions; additionally, we examine the role price connectedness plays in the setting of monetary policy. The third is \emph{methodological}: we replace the connectedness measures derived from the TVP-VAR with the
pseudo-$R^2$ quantile connectedness approach of \textcite{shahzadetal}, which better suits the
question for two reasons. First, it measures connectedness at different points of the
inflation distribution rather than only at the conditional mean or median, so the
dependence structure in high-inflation states can differ from that in normal
times.\footnote{A rolling-window TVP-VAR could in principle track the two episodes
through time, but it still measures connectedness at the centre of the conditional
distribution within each window; it cannot isolate behaviour in the tail.} Second, it
separates \emph{contemporaneous} co-movement from \emph{lagged} propagation, letting us
distinguish time series that move together within the period from those linked through
delayed pass-through; this is important in distinguishing common shocks from lagged propagation from price series to price series (that is, spiral type behaviour). Throughout, we emphasise that connectedness measures statistical
association and lead--lag dependence, not identified causal transmission; our claims
concern how the dependence structure shifts across inflation states, not what causes
what.

This paper tests six hypotheses. The first tests whether
inflation becomes more broad-based when inflation is high:

\noindent\textbf{H1 (CPI broadening).} As inflation rises, connectedness among
disaggregated CPI categories increases, particularly in the upper tail of the
inflation distribution.

The second asks whether high-inflation states are also associated with stronger
wage--price dependence:

\noindent\textbf{H2 (Wage--price interaction).} The connectedness between wages
and prices is stronger during high-inflation states than during normal-inflation
states.

We then compare two major U.S. inflation episodes: the Great Inflation of the
1970s and the COVID-era inflation surge of 2021--22. Much of the conventional
understanding of high-inflation dynamics is informed by the 1970s. The
pandemic-era surge provides a modern episode, under a very different
monetary policy regime, in which inflation became sufficiently high to revisit
those price/wage and inflation dynamics in high inflation environments. The two episodes differed in important respects. The Great
Inflation occurred in an economy more dependent on oil and a different monetary policy framework. The COVID-era inflation, differs from the Great Inflation, in that supply shock that affected many prices simultaneously; the pandemic-related supply constraints, in part owing to shipping disruptions was not just concentrated in oil-related phenomenon. These differences suggest that the two episodes may have
broadened through different channels. This motivates the third hypothesis:

\noindent\textbf{H3 (Episode mechanism).} The COVID-era surge was characterised
primarily by contemporaneous connectedness, consistent with a large common
shock, whereas the Great Inflation was characterised more by lagged propagation
across prices.

The differences in monetary policy frameworks between the two periods of high inflation motivate the fourth hypothesis. Specifically, we study if the adoption of inflation targeting and the associated anchoring of inflation expectations, alter the
connectedness of inflation expectations and inflation:

\noindent\textbf{H4 (Expectations channel).} Inflation expectations were a
stronger source of inflation propagation during the Great Inflation than during
the COVID-era surge.

Finally, we consider whether the breadth of price changes shapes the conduct of
monetary policy. As we noted earlier, there is an argument that monetary policy should
respond to broad-based, core inflation whilst `looking through' idiosyncratic
movements in individual components (relative price changes).\footnote{ See  \cite[p.276]{goodfriend1997new}} All else equal, then, when price changes are highly
connected across the subcomponents of the CPI, and as such are less likely to reflect idiosyncratic price changes, such
inflation may be harder to look through, and monetary policy makers may respond to it
more strongly:

\noindent\textbf{H5 (Connectedness and monetary policy).} When the subcomponents of the
CPI are more highly connected, monetary policy responds more strongly to inflation.

Further, when realised and expected inflation aree more tightly connected at high levels of inflation, the
risk of a self-reinforcing spiral rises: expectations may be behaving adaptively
rather than remaining anchored. Policy makers may then respond more aggressively
to forestall such a spiral. This motivates our final hypothesis:
 

\noindent\textbf{H6 (Connectedness and monetary policy II).} When inflation spills over
more strongly into inflation expectations at high levels of inflation, monetary policy responds more strongly to
inflation.


Our central finding is that the connectedness among prices changes with the
level of inflation. Cross-category connectedness among the CPI subcomponents rises
roughly threefold from the median to the high-inflation tail (from $11.8$ to $35.2$
per cent at $\tau=0.9$), and the increase is accounted for increasingly by the
\emph{lagged} rather than the contemporaneous component ($23.3$ against $11.9$ per
cent at $\tau=0.9$). We also find hetereogeneity in connectedness between the two periods of high inflation: the COVID-era surge chiefly
through contemporaneous co-movement, consistent with a common supply shock, and the
Great Inflation more through lagged propagation across categories. We find wages and prices are more tightly connected at the top end of the distribution, especially during the Great Inflation,  Further, as expected during the Great Inflation we see strong evidence consistent with a wage price spiral lagged inflation spills over into wages and lagged inflation spills over into inflation. We also find interesting results regarding the connectedness between inflation and inflation expectations.  Over the 1983--2019 period where it could be argued inflation expectations were well achonored, expectations transmit to inflation (net
connectedness $+29.1$ at $\tau=0.9$); but once  either high-inflation episode enters the sample period the sign reverses, so expectations \emph{follow}
realised inflation, consistent with adaptive rather than forward-looking behaviour. Finally, we some suggestive evidence that monetary policy responds more strongly to
inflation when CPI subcomponent price changes are more highly connected.  

The remainder of the paper proceeds as follows.
Section~\ref{sec:data} describes the data; Section~\ref{sec:method} sets out the method; and
Section~\ref{sec:results} reports the tests of the hypotheses.

% ============================================================
\section{Data}\label{sec:data}
% ============================================================

Hypotheses 1, 3 and 5 are to do with the connectedness between CPI subcomponents. Our study uses six monthly CPI subcomponent series; these series collectively  
cover all of the
headline CPI without overlapping. Each is expressed as a month-on-month percentage change and are downloaded from the FRED website. The sample covers
February 1967 to May 2026. The six subcomponents are:
\begin{itemize}
    \item Core goods (excluding food)
    \item Core services (excluding shelter)
    \item Shelter
    \item Food
    \item Household energy
    \item Gasoline
\end{itemize}

There are two exceptions from the standard expenditure-class decomposition: energy is separated into gasoline and household energy services, and core services into shelter and core services excluding shelter. The shelter series is available from FRED, but the core-services-ex-shelter series, we construct as:
\begin{align}
\%\Delta\text{CoreServ}^{\text{ex-shelter}}
  = \frac{\%\Delta\text{CoreServ} - w\,\%\Delta\text{Shelter}}{1-w},
\end{align}
where $w=0.60$ is estimated shelter's share of core services.\footnote{The results are insensitive to this choice: connectedness and the directional structure are stable for $w\in[0.50,0.65]$ (reported in the robustness analysis).}

We separate shelter from other core services because shelter inflation is driven by distinct dynamics and may behave differently from other services' prices.

To test hypotheses 2, 4 and 6, we use three additional series, again from FRED. For the overall Consumer
Price Index, we use the Consumer Price Index for All Urban Consumers, All Items, transformed into a
month-on-month percentage change to obtain headline inflation.\footnote{We omit the overall CPI from the subcomponent system because it is constructed from those very subcomponents, so any connectedness found would in part reflect a mechanical accounting relationship rather than economic co-movement.} Wage growth is measured using the average hourly earnings of production and nonsupervisory employees in the total private sector
(\texttt{AHETPI}, dollars per hour, seasonally adjusted), expressed as a
month-on-month percentage change. We use this series rather than the broader
all employees measure because it is the longest-running monthly earnings series,
beginning in 1964, and therefore spans the Great Inflation. A long span of data is needed for the
for comparative analysis of wage and inflation connectedness between the Great Inflation and the COVID Inflation. The broader
all-employees series begins only in 2006.

Finally, we use the University of Michigan median expected change in prices over the next twelve
months (\texttt{MICH}, percent), available monthly from January 1978. The monthly series therefore
misses the 1973--77 period, a key period of the Great Inflation. The Michigan measure is, however, available
back to 1960 at a quarterly frequency, so we also run the expectations analysis at a quarterly frequency.
   
    


% ============================================================
\section{Methodology}\label{sec:method}
% ============================================================

\subsection{Pseudo-quantile $R^2$ connectedness}


Let $\bm{y}_t=(y_{1,t},\dots,y_{K,t})'$ collect the $K$ time series we study. The precise membership of $\bm{y}_t$ depends on the hypothesis under study. For hypotheses 1, 3 and 5, $\bm{y}_t$ contains the month-on-month percentage change of the six CPI subcomponents. For hypotheses 4 and 5 it contains inflation expectations and headline inflation. For hypothesis 2 it contains wage growth  and headline inflation.

For each series $y_{k,t}$ in turn within the system $\bm{y}_t$, we ask how much of its variation can be attributed
to (i) variation of the remaining series in the system at time \textit{t}, $\bm{y}_{-k,t}$, and (ii)
the first $p$ lags of all the time series in the whole system, $(\bm{y}_{t-1}',\dots,\bm{y}_{t-p}')'$.
Channel (i) captures \emph{contemporaneous} dependence: the extent to which variables
move together within the month (or the quarter), while the second channel (ii) captures \emph{lagged}
dependence; that is, delayed propagation from one series to another over time. The ability to isolate and distinct between the two channels is what allows us to separate a common shock from
sequential pass-through.

A natural starting point to decompose the two channels is the goodness of fit of a regression as measured by the $R^2$. Consider quantifying the connectedness between inflation ($\pi$) and inflation expectations ($\pi^e$)
the following regression:

\[
\pi_t
=
\alpha_0
+
\alpha_1 \pi_t^e
+
\sum_{i=1}^{p} \beta_i \pi_{t-i}^e
%+
%\sum_{i=1}^{p} \gamma_i \pi_{t-i}
+
\epsilon_t.
\]

The $R^2$ tells us how well collectively  inflation expectations in time \textit{t} and lagged inflation expectations are correlated with inflation.  But the lags are correlated with the contemporaneous value of inflation expectations,  so the regression $R^2$ cannot be partitioned unambiguously into the portions that reflect the lagged and the contemporaneous contributions respectively.  This problem is exacerbated when you move to more than two variables (like the six individual subcomponents of the CPI) and therefore have contemporaneous and lagged values of two or more independent variables in the regression; if these variables are correlated, the partitioning of the $R^2$ becomes even more fraught.


\textcite{genizi1993} proposes an order-invariant decomposition of the regression
$R^2$ that distributes it across correlated regressors through the symmetric
(eigenvalue) square root of their correlation matrix. Let $\bm{R}_{xx}$ denote
the correlation matrix of the regressors, with eigendecomposition
$\bm{R}_{xx}=\bm{V}\bm{\Lambda}\bm{V}'$, where $\bm{V}$ holds the eigenvectors and
$\bm{\Lambda}$ the eigenvalues. Defining the symmetric square root $\bm{C}$, the
vector of decomposed $R^2$ shares is
\begin{align}
\bm{R}_{xx} &= \bm{V}\,\bm{\Lambda}\,\bm{V}' = \bm{C}\,\bm{C}', \label{eq:rxx}\\
\bm{C}      &= \bm{V}\,\bm{\Lambda}^{1/2}\,\bm{V}', \label{eq:C}\\
\bm{R}^{2,d} &= \bigl[\,\bm{C}^2\,(\bm{C}^{-1}\bm{r}_{yx})^2\,\bigr]' 
  \label{eq:r2d}
\end{align}
where $\bm{r}_{yx}$ is the vector of correlations between the dependent variable
and each regressor. Each entry of $\bm{R}^{2,d}$ is the share of the regression $R^2$ owed
uniquely to a single regressor, net of the correlation it shares with the others; the entries are
non-negative and sum to the regression $R^2$.

%\subsubsection{Measuring dependence in the tails}
 Using the above decomposition with Pearson correlations as inputs, \textcite{balli2023} show how contemporaneous and lagged connectedness can be quantified at the centre of the distribution. Because, however, our interest is in how dependence behaves when inflation is high, we use the adaptation of \textcite{shahzadetal}, which replaces the Pearson correlations with the quantile correlations of \textcite{choishin2022}; this allows contemporaneous and lagged connectedness to be measured at any chosen quantile $\tau$.

The specifics of the \textcite{shahzadetal} measure are as follows. The quantile correlation of \textcite{choishin2022} between two variables, $r_{1,2,\tau}$, is substituted for $\bm{R}_{xx}$ and $\bm{r}_{yx}$ in equations~\ref{eq:rxx} and~\ref{eq:r2d}. For each time series in turn,
equation~\ref{eq:r2d} then returns a vector $\bm{R}^{2,d}_{\tau}$ that allocates
the regression's overall pseudo-$R^2$ across its regressors at that $\tau$.\footnote{Quantile regressions do \emph{not} have an $R^2$ in the conventional (least-squares) sense, hence the term ``pseudo''.} The first $K-1$ entries
record the contemporaneous contributions of the other series; the remaining
entries record the lagged contributions of all series in the system.


For a system of $K$ variables, applying  the decomposition with the quantile correlation for each of the $K$ equations in that system and stacking the results
yields the connectedness matrix
\begin{align}
\bm{R}_{\tau}^{2,d}
  = \begin{bmatrix}
      \bm{R}_{\tau,1}^{2,d} \\ \vdots \\ \bm{R}_{\tau,K}^{2,d}
    \end{bmatrix}
  = \bigl[\, \bm{R}_{lag=0,\tau}^{2,d} \;\big|\; \bm{R}_{lag=1,\tau}^{2,d}
    \;\big|\; \cdots \;\big|\; \bm{R}_{lag=p,\tau}^{2,d} \,\bigr],
\end{align}
in which row $k$ indexes the dependent variable.\footnote{So that the equations align, the $k$th position
of $\bm{R}_{\tau,k}^{2,d}$ is set to zero, since $y_{k,t}$ is not regressed on its
own contemporaneous value.}   


The total connectedness index at quantile $\tau$, $\mathrm{TCI}_{\tau}$, is the average across the
$K$ series of the off-diagonal $R^{2,d}$ shares, and splits into a contemporaneous and a lagged
component,
\begin{align}
\mathrm{TCI}_{\tau} = \mathrm{TCI}_{\tau}^{C} + \mathrm{TCI}_{\tau}^{L},
\end{align}
with
\begin{align}
\mathrm{TCI}_{\tau}^{C}
  = \frac{1}{K}\sum_{k=1}^{K}\sum_{j\neq k}\bm{R}_{lag=0,\tau,k,j}^{2,d},
\qquad
\mathrm{TCI}_{\tau}^{L}
  = \frac{1}{K}\sum_{k=1}^{K}\sum_{j=1}^{K}\sum_{l=1}^{p}
    \bm{R}_{lag=l,\tau,k,j}^{2,d},
\end{align}
where $\bm{R}_{lag=l,\tau,k,j}^{2,d}$ is the share of series $k$'s variation contributed by
series $j$ at lag $l$. Throughout, connectedness measures are reported as percentages (shares of
the $R^2$ multiplied by 100). Equivalently, $\mathrm{TCI}_{\tau}$ is the average of the directional
\textsc{from} --- or, identically, \textsc{to} --- measures defined below [in this true?].

Directional measures in a similar spirit to those put forward by \textcite{dieboldyilmaz2012} can also be calculated. The total connectedness \emph{received} by series $k$
(its \textsc{from} measure) sums the shares contributed by all other series, at
every lag, into row $k$; 
\begin{align}
\mathrm{FROM}_{k,\tau}
  &= \sum_{j\neq k}\bm{R}_{lag=0,\tau,k,j}^{2,d}
   + \sum_{j=1}^{K}\sum_{l=1}^{p}\bm{R}_{lag=l,\tau,k,j}^{2,d}.
\end{align}

[now we switch to j instead of k, would we? ]The total connectedness \emph{transmitted} by series $j$ to other series
(its \textsc{to} measure) sums the shares it contributes, at every lag, down
column $j$:
\begin{align}
\mathrm{TO}_{j,\tau}
  &= \sum_{k\neq j}\bm{R}_{lag=0,\tau,k,j}^{2,d}
   + \sum_{k=1}^{K}\sum_{l=1}^{p}\bm{R}_{lag=l,\tau,k,j}^{2,d}.
\end{align}

The net directional connectedness of series $j$,
\begin{align}
\mathrm{NET}_{j,\tau} = \mathrm{TO}_{j,\tau} - \mathrm{FROM}_{j,\tau},
\end{align}
is positive when $j$ is, on balance, a transmitter of connectedness and
negative when it is a net receiver.

We prefer this measure of  quantile connectednes to other approaches to quantile connectedness, such as that of \textcite{ando2022}. Those measures propagate the non-zero-mean residuals from quantile regressions through a generalised forecast-error variance decomposition that treats them as zero-mean shocks. As \textcite{shahzadetal} show, this inflates measured connectedness in both tails.


\section{Results}\label{sec:results}

All connectedness measures reported below are percentages (shares of the $R^2$, multiplied by
$100$). 

\subsection{H1: inflation broadens in the high tail}
Our first hypothesis is \textbf{H1 (CPI broadening)}: as inflation rises,
connectedness among disaggregated CPI categories increases, particularly in the
upper tail of the inflation distribution. Table~\ref{tab:h1} reports the overall
connectedness at different levels of $\tau$ over our whole sample period (the
and its decomposition into contemporaneous and lagged 
connectedness; Table~\ref{tab:h1} reports the same by quantile, monthly (Panel~A)
and quarterly (Panel~B). Overall connectedness rises about threefold from the
median to the high tail (from $11.8$ to $35.2$ at monthly frequency); both the
contemporaneous and lagged components increase, but slightly more of the rise
comes from the lagged component ($23.3$ vs $11.9$ at $\tau=0.9$). This indicates that high-inflation states are characterised not only by stronger common movements across prices, but also by more pronounced sequential propagation across categories. This result also holds if quarterly data are used (see Panel B).

\begin{figure}[!ht]
\centering
\includegraphics[width=0.68\linewidth]{paper/TCI_decomp.pdf}
\caption{Overall connectedness (black) and its decomposition into
contemporaneous (blue) and lagged (red) components, by inflation quantile
$\tau$ (monthly, $n_{\text{lag}}=2$).}
\label{fig:h1}
\end{figure}

Table \ref{conn_contemp_full}    in the appendix show the contemporaneous dependency between CPI subcomponents at $\tau = 0.9$. The rows indicate FROM connectedness so Food gets 0.2 units from Gas and 1.6 from Energy; the columns indicate TO connectedness, so Food gives 0.4 to gasoline and 1.6 household energy. The strongest contemporaneous connectedness is recorded between household energy and gasoline (in both directions), and  core goods and core services with shelter, again in both directions.  The bidirectional nature of these suggests they represent common shocks.  Higher nominal incomes and low interest rates could increase demand for houses, good and services for example.     
Table [] shows the same for lagged connectedness. In line with the results in Table~\ref{tab:h1} we see, generally, connectedness is stronger with a lag.  WE see that shelter, services and gasoline are a net transmitter of non-contemporaneous connectedness, where as Food, core goods and household energy are net receivers [need to revise this].

\begin{table}[!ht]
\centering
\caption{H1 --- Inflation broadens into the high tail}
\label{tab:h1}
\setlength{\tabcolsep}{6pt}\footnotesize
\begin{threeparttable}
\begin{tabular}{c c c c}
\multicolumn{4}{l}{\textit{Panel A. Monthly total connectedness by quantile}}\\
\toprule
$\tau$ & Overall & Contemporaneous & Lagged \\
\midrule
0.1 & 5.5 & 2.8 & 2.8 \\
0.3 & 7.5 & 2.7 & 4.8 \\
0.5 & 11.8 & 4.7 & 7.2 \\
0.7 & 19.7 & 7.1 & 12.6 \\
0.9 & 35.2 & 11.9 & 23.3 \\
\bottomrule
\end{tabular}

\vspace{4pt}
\begin{tabular}{c c c c}
\multicolumn{4}{l}{\textit{Panel B. Quarterly total connectedness by quantile}}\\
\toprule
$\tau$ & Overall & Contemporaneous & Lagged \\
\midrule
0.1 & 16.9 & 8.3 & 8.6 \\
0.3 & 17.6 & 6.8 & 10.8 \\
0.5 & 27.7 & 10.7 & 17.1 \\
0.7 & 55.7 & 21.6 & 34.1 \\
0.9 & 86.8 & 33.8 & 52.9 \\
\bottomrule
\end{tabular}
\begin{tablenotes}[flushleft]\footnotesize
\item \textit{Notes:} Pseudo-quantile $R^2$ (Genizi/nearPD) connectedness (\%) among the six CPI subcomponents; Overall $=$ Contemporaneous $+$ Lagged, $n_{\text{lag}}=2$. Panel A is monthly 1967--2026 (broadening ratio $\text{TCI}(0.9)/\text{TCI}(0.5)=2.98\times$); Panel B quarterly 1967Q2--2026Q2 (ratio 3.13$\times$). Both frequencies show the tail rise, increasingly driven by the lagged component. Point estimates; bootstrap confidence intervals for the six-CPI system are to be added.
\end{tablenotes}
\end{threeparttable}
\end{table}



\subsection{H3: how the two episodes broadened}
Our most distinctive result concerns \emph{how} inflation broadened in each episode.
Table~\ref{tab:h3} reports how high-tail ($\tau=0.9$) connectedness changes when each episode is
added to the calm 1983--2019 core --- estimating the model with the Great Inflation included (but
COVID dropped) and with COVID included (but the Great Inflation dropped). Both contemporaneous and
lagged connectedness rise by more in the Great Inflation, consistent with its higher inflation
level. More tellingly, the Great-Inflation increase comes mainly through the \emph{lagged}
component, whereas the COVID increase is almost entirely \emph{contemporaneous} --- its lagged
marginal is in fact negative. The two components have a natural reading: contemporaneous
connectedness is categories moving together within the period, as a common shock hitting all prices
at once would produce; lagged connectedness is one category's past movement being associated with
another's current movement, the signature of a sequential, cross-category process. The COVID pattern
is thus consistent with a common shock (pandemic supply constraints and shipping disruption), and
the Great-Inflation pattern with a slower, sequential build-up --- supporting H3. This distinction
is precisely the object the contemporaneous/lagged decomposition is built to measure, and it links
to the wage channel examined next. The contrast is sharp at monthly frequency (Panel~A); at
quarterly frequency (Panel~B) the calm-era base connectedness is already high and lagged-dominated,
so the episode marginals are small and less sharply signed, and we read the monthly result as the
primary evidence.

\begin{table}[!ht]
\centering
\caption{H3 --- COVID broadened contemporaneously; the Great Inflation through lagged propagation}
\label{tab:h3}
\setlength{\tabcolsep}{6pt}\footnotesize
\begin{threeparttable}
\begin{tabular}{l c c c c}
\multicolumn{5}{l}{\textit{Panel A. Monthly: marginal $\Delta$TCI$(\tau)$ vs.\ the 1983--2019 core}}\\
\toprule
Episode & $\Delta$Overall & $\Delta$Contemp. & $\Delta$Lagged & Lagged \% \\
\midrule
\multicolumn{5}{l}{\textit{\;$\tau=0.7$}}\\
GreatInflation (1967-82) & +12.3 & +4.4 & +7.9 & 64.4 \\
COVID (2020-26) & -0.1 & +0.0 & -0.2 & 139.6 \\
\addlinespace
\multicolumn{5}{l}{\textit{\;$\tau=0.9$}}\\
GreatInflation (1967-82) & +14.3 & +5.0 & +9.3 & 64.9 \\
COVID (2020-26) & +1.1 & +1.9 & -0.8 & -74.7 \\
\bottomrule
\end{tabular}

\vspace{4pt}
\begin{tabular}{l c c c c}
\multicolumn{5}{l}{\textit{Panel B. Quarterly: marginal $\Delta$TCI$(\tau)$ vs.\ the 1983--2019 core}}\\
\toprule
Episode & $\Delta$Overall & $\Delta$Contemp. & $\Delta$Lagged & Lagged \% \\
\midrule
\multicolumn{5}{l}{\textit{\;$\tau=0.7$}}\\
GreatInflation (1967-82) & +25.6 & +6.9 & +18.6 & 72.9 \\
COVID (2020-26) & -0.1 & -0.4 & +0.3 & -405.9 \\
\addlinespace
\multicolumn{5}{l}{\textit{\;$\tau=0.9$}}\\
GreatInflation (1967-82) & +3.2 & +3.8 & -0.6 & -17.8 \\
COVID (2020-26) & -19.9 & -5.0 & -14.9 & 74.9 \\
\bottomrule
\end{tabular}
\begin{tablenotes}[flushleft]\footnotesize
\item \textit{Notes:} Each episode's marginal contribution to connectedness at $\tau\in\{0.7,0.9\}$ relative to the calm 1983--2019 core (nested subtraction), split into contemporaneous and lagged; six-CPI system, Panel A monthly, Panel B quarterly. At monthly frequency the Great Inflation broadened mainly through the \emph{lagged} (propagation) channel at both quantiles, while COVID broadened contemporaneously (its lagged marginal is negative at $\tau=0.9$ and essentially zero at $\tau=0.7$). At quarterly frequency the same contrast is sharp at $\tau=0.7$; at $\tau=0.9$ the calm-era base is already high and lagged-dominated, and the COVID window contributes only 26 quarters, so the episode marginals there are less informative.
\end{tablenotes}
\end{threeparttable}
\end{table}



\subsection{Self-reinforcement: wages and expectations}
The next two hypotheses ask whether inflation became self-reinforcing through wages and
expectations. Table~\ref{tab:h2} shows that the connectedness between wages and the overall CPI
rises with the quantile (H2), and that at higher quantiles the lagged inflation--wage association
strengthens. The pattern holds at quarterly frequency (Panel~C) and is concentrated in the
Great-Inflation-inclusive subperiod (Panel~B): wage--price connectedness at $\tau=0.9$ is $25.2$ in
1964--2019 against only $3.6$ post-1983. Wage--price dependence is thus, like broadening itself, a
high-inflation-regime phenomenon.

\begin{table}[!ht]
\centering
\caption{H2 --- Wage--price connectedness is stronger when inflation is high}
\label{tab:h2}
\setlength{\tabcolsep}{6pt}\footnotesize
\begin{threeparttable}
\begin{tabular}{l c c c}
\multicolumn{4}{l}{\textit{Panel A. Two-variable wage--CPI system (monthly, full sample), by quantile}}\\
\toprule
 & $\tau=0.5$ & $\tau=0.7$ & $\tau=0.9$ \\
\midrule
Wage--CPI total connectedness (TCI) & 10.8 & 16.6 & 21.6 \\
\quad lagged Wages$\to$CPI & 5.0 & 7.7 & 11.5 \\
\quad lagged CPI$\to$Wages & 6.2 & 9.2 & 16.1 \\
\bottomrule
\end{tabular}

\vspace{4pt}
\begin{tabular}{l c c c}
\multicolumn{4}{l}{\textit{Panel B. By subperiod (monthly)}}\\
\toprule
 & $\tau=0.5$ & $\tau=0.7$ & $\tau=0.9$ \\
\midrule
\multicolumn{4}{l}{\textit{\;1964--2019 (with Great Inflation)}}\\
Wage--CPI total connectedness (TCI) & 13.3 & 18.4 & 25.2 \\
\quad lagged Wages$\to$CPI & 6.0 & 9.1 & 15.5 \\
\quad lagged CPI$\to$Wages & 5.9 & 9.0 & 14.9 \\
\addlinespace
\multicolumn{4}{l}{\textit{\;1983--2026 (post-Great-Inflation)}}\\
Wage--CPI total connectedness (TCI) & 1.2 & 1.8 & 3.6 \\
\quad lagged Wages$\to$CPI & 0.2 & 0.2 & 1.3 \\
\quad lagged CPI$\to$Wages & 0.9 & 1.7 & 3.0 \\
\bottomrule
\end{tabular}

\vspace{4pt}
\begin{tabular}{l c c c}
\multicolumn{4}{l}{\textit{Panel C. Quarterly (full sample), by quantile}}\\
\toprule
 & $\tau=0.5$ & $\tau=0.7$ & $\tau=0.9$ \\
\midrule
Wage--CPI total connectedness (TCI) & 25.7 & 32.6 & 53.4 \\
\quad lagged Wages$\to$CPI & 14.6 & 19.3 & 35.1 \\
\quad lagged CPI$\to$Wages & 13.0 & 18.5 & 23.8 \\
\bottomrule
\end{tabular}
\begin{tablenotes}[flushleft]\footnotesize
\item \textit{Notes:} Two-variable system of wage growth (AHETPI) and overall CPI (CPIAUCSL); pseudo-quantile $R^2$ connectedness (\%), Genizi estimator, $n_{\text{lag}}=2$. Wage--price connectedness rises sharply into the high-inflation tail at all frequencies, and is far stronger in the Great-Inflation-inclusive subperiod (1964--2019) than post-1983. In the high tail the lagged components dominate. Quarterly ($n_{\text{lag}}=2$) spans a six-month horizon.
\end{tablenotes}
\end{threeparttable}
\end{table}



For expectations (\textbf{H4}), Table~\ref{tab:h6net} reports the net directional connectedness of
inflation expectations. Pooled over the full sample, expectations are net transmitters in the upper
tail --- more connected \emph{to} inflation than they receive \emph{from} it --- but this reflects
the anchored 1983--2019 period, where expectations lead (NET $=+29.1$ at $\tau=0.9$). Conditional on
a high-inflation episode the sign reverses: in the high tail of both the Great Inflation (NET
$=-20.2$) and COVID ($-6.7$) expectations follow realised inflation, consistent with adaptive
expectations. The Great-Inflation figure uses the quarterly Michigan series (from 1960), which
covers the 1973--77 un-anchoring the monthly series misses. Table~\ref{tab:h6} decomposes the link
by direction: in the full sample the lagged expectations$\to$inflation term is the larger, but in
the Great Inflation the lagged inflation$\to$expectations term dominates at higher $\tau$ (reaching
$50.7$ at $\tau=0.9$, against $27.5$ in the reverse direction) --- the adaptive channel. COVID shows
the same pattern, less pronounced.

\begin{table}[!ht]
\centering
\caption{H4 --- Do inflation expectations lead or follow? Net directional connectedness}
\label{tab:h6net}
\setlength{\tabcolsep}{6pt}\footnotesize
\begin{threeparttable}
\begin{tabular}{l c c c}
\toprule
Sample & $\tau=0.5$ & $\tau=0.7$ & $\tau=0.9$ \\
\midrule
Full (1960--2026) & +4.8 & +13.7 & +10.5 \\
Great Inflation (1967--82) & +0.5 & -2.2 & -20.2 \\
Core (1983--2019) & +3.0 & +8.8 & +29.1 \\
COVID (2020--26) & -5.5 & -4.1 & -6.7 \\
\bottomrule
\end{tabular}
\begin{tablenotes}[flushleft]\footnotesize
\item \textit{Notes:} Net directional connectedness of quarterly Michigan inflation expectations (\texttt{MICH\_QTR}, from 1960) versus overall CPI in a two-variable system, $n_{\text{lag}}=2$. $\text{NET}=\text{TO}-\text{FROM}$; $>0$ means expectations lead (transmit), $<0$ means they follow. Pooled, expectations lead in the upper tail --- but this reflects the anchored 1983--2019 core (NET $+29.1$ at $\tau=0.9$); conditional on a high-inflation episode the sign reverses (Great Inflation $-20.2$, COVID $-6.7$), consistent with adaptive expectations. The COVID window (25 quarters) is short.
\end{tablenotes}
\end{threeparttable}
\end{table}


\begin{table}[!ht]
\centering
\caption{H4 --- Contemporaneous and lagged connectedness between inflation and expectations}
\label{tab:h6}
\setlength{\tabcolsep}{6pt}\footnotesize
\begin{threeparttable}
\begin{tabular}{l c c c c c c}
\toprule
 & & \multicolumn{2}{c}{Inflation $\to$ Expectations} & \multicolumn{2}{c}{Expectations $\to$ Inflation} & \\
\cmidrule(lr){3-4}\cmidrule(lr){5-6}
Sample & $\tau$ & Contemp. & Lagged & Contemp. & Lagged & NET \\
\midrule
Full (1960--2026) & 0.5 & 24.2 & 23.7 & 27.4 & 25.4 & +4.8 \\
 & 0.7 & 20.9 & 28.5 & 27.7 & 35.4 & +13.7 \\
 & 0.9 & 19.5 & 33.6 & 26.2 & 37.3 & +10.5 \\
 & 0.78$^{\ddagger}$ & 20.2 & 28.7 & 28.0 & 37.9 & +16.9 \\
\addlinespace
Core (1983--2019) & 0.5 & 17.9 & 8.8 & 22.8 & 6.8 & +3.0 \\
 & 0.7 & 14.9 & 6.2 & 21.6 & 8.2 & +8.8 \\
 & 0.9 & 21.3 & 11.9 & 27.5 & 34.7 & +29.1 \\
 & 0.95$^{\ddagger}$ & 21.5 & 20.5 & 25.2 & 24.2 & +7.3 \\
\addlinespace
1960--2019 & 0.5 & 24.3 & 26.8 & 26.4 & 28.6 & +3.9 \\
 & 0.7 & 21.0 & 28.8 & 25.2 & 33.0 & +8.3 \\
 & 0.9 & 20.9 & 36.6 & 26.5 & 40.3 & +9.3 \\
 & 0.79$^{\ddagger}$ & 20.9 & 30.6 & 28.1 & 36.9 & +13.6 \\
\addlinespace
1970--2019 & 0.5 & 25.3 & 23.0 & 27.3 & 26.8 & +5.8 \\
 & 0.7 & 20.5 & 28.3 & 24.6 & 32.1 & +8.0 \\
 & 0.9 & 19.0 & 40.7 & 27.5 & 41.1 & +8.9 \\
 & 0.77$^{\ddagger}$ & 20.5 & 30.6 & 27.2 & 37.2 & +13.3 \\
\addlinespace
1983--2026 & 0.5 & 15.4 & 8.4 & 21.5 & 7.9 & +5.7 \\
 & 0.7 & 11.4 & 14.6 & 16.1 & 8.6 & -1.4 \\
 & 0.9 & 18.0 & 16.1 & 38.0 & 47.8 & +51.8 \\
 & 0.92$^{\ddagger}$ & 23.7 & 13.5 & 42.6 & 42.2 & +47.6 \\
\addlinespace
1993--2026 & 0.5 & 7.1 & 6.3 & 9.4 & 3.1 & -0.9 \\
 & 0.7 & 9.2 & 13.5 & 14.5 & 12.5 & +4.4 \\
 & 0.9 & 25.3 & 16.3 & 45.3 & 43.2 & +47.0 \\
 & 0.93$^{\ddagger}$ & 16.6 & 20.9 & 45.2 & 37.8 & +45.5 \\
\bottomrule
\end{tabular}
\begin{tablenotes}[flushleft]\footnotesize
\item \textit{Notes:} Pairwise pseudo-quantile $R^2$ connectedness (\%) in a two-variable quarterly system of Michigan inflation expectations (\texttt{MICH\_QTR}, from 1960) and overall CPI (CPIAUCSL), $n_{\text{lag}}=2$ (a six-month lag horizon). ``Inflation $\to$ Expectations'' is the share of \emph{expectations'} variation explained by CPI, split into contemporaneous and lagged; ``Expectations $\to$ Inflation'' is the reverse. NET is the net directional connectedness of expectations (TO $-$ FROM); $>0$ means expectations lead (transmit), $<0$ that they follow. Dependence, not identified causation. $^{\ddagger}$~The fourth row of each block reports connectedness at the quantile whose quarter-on-quarter CPI inflation equals $5$ per cent annualised within that sample (a common-level comparison across samples). $^{\dagger}$~The $5$ per cent level lies outside $[0.05,0.95]$ for that sample and the quantile is clamped to the boundary.
\end{tablenotes}
\end{threeparttable}
\end{table}




To test H5 and H6 we proceed in three steps. First, for each connectedness system we
compute the pseudo-quantile $R^2$ connectedness measure over rolling windows, using both
60- and 120-month windows and evaluating it at $\tau=0.5$, $0.7$ and $0.9$ so as to
characterise connectedness at the centre as well as in the upper tail of the distribution.
This yields, for each window length and quantile, a monthly series of total connectedness,
$C_t$.

Second, we classify each month into a high- or low-connectedness regime. Rather than
imposing an ad hoc threshold, we let the data determine the regimes by fitting the
two-state Markov-switching model of \citet{hamilton1989new} to $C_t$. Specifically, we
treat connectedness as a first-order autoregression whose intercept, autoregressive
coefficient and error variance all switch between two latent states,
$s_t\in\{\text{low},\text{high}\}$, that evolve as a first-order Markov chain. The AR(1)
specification captures the well-known persistence of the connectedness series whilst
allowing both its level and its volatility to differ across regimes. We estimate the model
by maximum likelihood, recover the smoothed probability that each month belongs to each
state, and assign every month to its most probable state; the state with the higher mean
connectedness is labelled the `high' regime, so that $\text{High}_t=1$ for months in that
state and zero otherwise.

Third,   we estimate an inertial Taylor rule augmented
with a regime interaction:
\begin{equation}
  i_t = \alpha + \rho\,i_{t-1} + \beta_\pi\,(\pi_t-\pi^*)
        + \delta\,\big(\text{High}_t\cdot(\pi_t-\pi^*)\big) + \beta_x\,x_t + \varepsilon_t,
  \label{eq:regime_taylor}
\end{equation}
where $i_t$ is the Wu--Xia \citep{wu2016measuring} shadow policy rate, $\pi_t$ is twelve-month
CPI inflation, $\pi^*=2$ per cent is the assumed inflation target, and $x_t$ is the
unemployment gap.\footnote{he unemployment gap is the difference between the civilian unemployment rate
and the Congressional Budget Office's estimate of the natural (non-cyclical) rate of
unemployment, both obtained from FRED (series \texttt{UNRATE} and \texttt{NROU}). Because
the CBO series is published quarterly, we linearly interpolate it to monthly frequency.} The coefficient $\rho$ captures interest-rate smoothing  \cite[p.151-152]{ClaridaGaliGertler2000}; one motivation for this specification is that under uncertainty central banks only partially adjust to the interest rate warranted based on current economic conditions. In equation   \ref{eq:regime_taylor} $\beta_\pi$ the
baseline response to the inflation gap, and $\beta_x$ the response to unemployment. The
object of interest is $\delta$: it measures the \emph{additional} response to inflation in
the high-connectedness regime. If H5 and H6 are true, then $\delta>0$: monetary policy
responds more strongly to inflation when price changes (H5) or inflation and inflation expectations (H6) are more highly connected. All inference is based on
Newey--West HAC standard errors.  

 

Table~\ref{tab:sixconn} reports estimates of equation~\eqref{eq:regime_taylor} in which
the high- and low-connectedness regimes are identified from the connectedness of the six
CPI subcomponents, measured over a 120-month rolling window. We re-estimate the rule over
several sample start dates because the monetary-policy reaction function is known to have
shifted over time, with different Federal Reserve chairs placing different weight on
inflation \citep{ClaridaGaliGertler2000}; estimating over later subsamples guards against
attributing such governor-preference changes in the rule to the connectedness regime.  

 



Two readings of this evidence are possible. On the first, the state dependence is rea;
but recent: it emerges only in the samples beginning in 1993 and 2003, consistent with
modern central banks, operating under explicit or implicit inflation targets. being more
attentive to the breadth of price changes than their predecessors. That the effect appears
only at $\tau=0.9$ reinforces this reading, suggesting that policymakers respond more
strongly to inflation only once connectedness is concentrated in the upper tail, amongst
those components already experiencing large price increases. This interpretation has some plausiblity. A second, more sceptical reading is that the two significant coefficients
are false positives. The interaction $\delta$ is estimated over twelve quantile--sample
combinations, so at conventional levels one or two spurious rejections are to be expected
by chance alone; that one of the two is significant only at the 10 per cent level lends
this reading further weight. We therefore treat the result as no more than suggestive, and
note below that it does not survive a shorter, 60-month estimation window.

Table \label{tab:expconn} shows connectedness between inflation expectations and inflation under two different measures of inflation.  One the annual percentage change in inflation (panel A), one is the monthly change in inflatioon (panel B); we study the annual percentage change so it is one the same period as inflation expectations which are measured over a year.  On that measure only $\tau=0.7$ connectedness shows any change in the monetary policy reaction function for inflation; this does not survive if monthly CPI growth is used. 

Table~\ref{tab:expconn} reports the Taylor rule estimates when the connectedness
regime is identified from the inflation--expectations system, under two measures of
inflation: the twelve-month (annual) percentage change in the CPI (Panel~A) and the
month-on-month change (Panel~B). We prefer the annual measure, as it is defined over the
same twelve-month horizon as the Michigan expectations series. On that measure the
interaction $\delta$ is positive and statistically significant only at $\tau=0.7$, where
it is significant across every sample start date, with $\delta$ ranging from $0.04$ to
$0.06$ indicating a stronger inflation response in high-connectedness months. This
result does not survive when month-on-month CPI inflation is used instead: in Panel~B
$\delta$ is statistically insignificant throughout. Again not overly strong support for the hypothesis that connectedness is taken into account when setting monetary policy.

\begin{table}[!htbp]
\centering
\caption{Subcomponent connectedness }
\label{tab:expconn}

\setlength{\tabcolsep}{4pt}\scriptsize
\begin{tabular}{lcccc}
\toprule
 & Earliest (1977) & 1983 & 1993 & 2003 \\
\midrule
\multicolumn{5}{l}{\textit{$\tau=0.5$}} \\
Intercept & $0.16^{**}$\,{\scriptsize(2.1)} & $0.10^{*}$\,{\scriptsize(1.7)} & $0.04$\,{\scriptsize(0.9)} & $0.04$\,{\scriptsize(1.1)} \\
$\rho$ ($i_{t-1}$) & $0.95^{***}$\,{\scriptsize(54.3)} & $0.96^{***}$\,{\scriptsize(62.3)} & $0.99^{***}$\,{\scriptsize(96.5)} & $0.98^{***}$\,{\scriptsize(95.4)} \\
$\beta_\pi$ ($\pi-\pi^*$) & $0.06^{**}$\,{\scriptsize(2.3)} & $0.02$\,{\scriptsize(1.0)} & $-0.03$\,{\scriptsize(-0.6)} & $0.02$\,{\scriptsize(1.1)} \\
$\Delta\beta_\pi^{\text{High}}$ & $0.01$\,{\scriptsize(0.3)} & $0.03$\,{\scriptsize(1.4)} & $0.07$\,{\scriptsize(1.3)} & $0.03$\,{\scriptsize(1.4)} \\
$\beta_x$ (ugap) & $-0.04^{**}$\,{\scriptsize(-2.4)} & $-0.03^{*}$\,{\scriptsize(-1.8)} & $-0.02$\,{\scriptsize(-1.5)} & $-0.03^{**}$\,{\scriptsize(-2.3)} \\
\addlinespace
$n$ & 590 & 518 & 398 & 278 \\
high-conn.\ share & 0.40 & 0.33 & 0.63 & 0.15 \\
long-run $\beta_\pi$ (low/high) & $1.19/1.31$ & $0.53/1.41$ & $-2.91/3.92$ & $0.94/2.52$ \\
\midrule
\multicolumn{5}{l}{\textit{$\tau=0.7$}} \\
Intercept & $0.16^{**}$\,{\scriptsize(2.2)} & $0.11^{*}$\,{\scriptsize(1.8)} & $0.04$\,{\scriptsize(0.9)} & $0.04$\,{\scriptsize(1.0)} \\
$\rho$ ($i_{t-1}$) & $0.95^{***}$\,{\scriptsize(54.7)} & $0.97^{***}$\,{\scriptsize(61.6)} & $0.99^{***}$\,{\scriptsize(98.4)} & $0.98^{***}$\,{\scriptsize(97.5)} \\
$\beta_\pi$ ($\pi-\pi^*$) & $0.03$\,{\scriptsize(1.1)} & $-0.01$\,{\scriptsize(-0.4)} & $-0.02$\,{\scriptsize(-0.5)} & $0.02$\,{\scriptsize(1.3)} \\
$\Delta\beta_\pi^{\text{High}}$ & $0.04$\,{\scriptsize(1.5)} & $0.06$\,{\scriptsize(1.6)} & $0.06$\,{\scriptsize(1.2)} & $0.03$\,{\scriptsize(1.3)} \\
$\beta_x$ (ugap) & $-0.04^{**}$\,{\scriptsize(-2.5)} & $-0.03^{*}$\,{\scriptsize(-1.8)} & $-0.02$\,{\scriptsize(-1.5)} & $-0.03^{**}$\,{\scriptsize(-2.2)} \\
\addlinespace
$n$ & 590 & 518 & 398 & 278 \\
high-conn.\ share & 0.62 & 0.62 & 0.58 & 0.16 \\
long-run $\beta_\pi$ (low/high) & $0.59/1.31$ & $-0.41/1.29$ & $-2.20/3.67$ & $1.34/2.81$ \\
\midrule
\multicolumn{5}{l}{\textit{$\tau=0.9$}} \\
Intercept & $0.17^{**}$\,{\scriptsize(2.2)} & $0.10$\,{\scriptsize(1.6)} & $0.03$\,{\scriptsize(0.9)} & $0.04$\,{\scriptsize(1.3)} \\
$\rho$ ($i_{t-1}$) & $0.95^{***}$\,{\scriptsize(55.4)} & $0.97^{***}$\,{\scriptsize(61.9)} & $0.99^{***}$\,{\scriptsize(116.0)} & $0.98^{***}$\,{\scriptsize(104.4)} \\
$\beta_\pi$ ($\pi-\pi^*$) & $0.04$\,{\scriptsize(1.2)} & $0.03$\,{\scriptsize(1.2)} & $0.01$\,{\scriptsize(0.4)} & $0.01$\,{\scriptsize(0.7)} \\
$\Delta\beta_\pi^{\text{High}}$ & $0.03$\,{\scriptsize(1.1)} & $0.02$\,{\scriptsize(0.6)} & $0.04^{**}$\,{\scriptsize(2.0)} & $0.04^{*}$\,{\scriptsize(1.8)} \\
$\beta_x$ (ugap) & $-0.04^{**}$\,{\scriptsize(-2.3)} & $-0.03^{*}$\,{\scriptsize(-1.7)} & $-0.02^{*}$\,{\scriptsize(-1.8)} & $-0.03^{**}$\,{\scriptsize(-2.6)} \\
\addlinespace
$n$ & 590 & 518 & 398 & 278 \\
high-conn.\ share & 0.45 & 0.37 & 0.26 & 0.31 \\
long-run $\beta_\pi$ (low/high) & $0.67/1.33$ & $0.82/1.32$ & $0.45/3.52$ & $0.57/2.42$ \\
\bottomrule
\end{tabular}

\begin{tablenotes}\footnotesize
\item \textit{Notes:} Estimates of equation~\eqref{eq:regime_taylor}. Columns are sample
start dates; each panel a quantile $\tau$. Newey--West HAC $t$-statistics in parentheses;
$^{*}$, $^{**}$, $^{***}$ denote significance at 10, 5 and 1 per cent. $\Delta\beta_\pi^{
\text{High}}$ is the interaction $\delta$; the final row reports the implied long-run
response $\beta_\pi/(1-\rho)$ in the low and high regimes.
\end{tablenotes}
\end{table}

\begin{table}[!htbp]
\centering
\caption{Expectations connectedness and the policy reaction function: 120-month window}
\label{tab:expconn120}
\begin{threeparttable}
\centering

\textbf{Panel A. Matched: MICH and 12-month CPI inflation}\par\smallskip
{\setlength{\tabcolsep}{4pt}\scriptsize
\begin{tabular}{lccc}
\toprule
 & Earliest (1987) & 1993 & 2003 \\
\midrule
\multicolumn{4}{l}{\textit{$\tau=0.5$}} \\
Intercept & $0.05$\,{\scriptsize(1.4)} & $0.03$\,{\scriptsize(0.9)} & $0.04$\,{\scriptsize(1.3)} \\
$\rho$ ($i_{t-1}$) & $0.98^{***}$\,{\scriptsize(112.3)} & $0.98^{***}$\,{\scriptsize(111.3)} & $0.98^{***}$\,{\scriptsize(99.2)} \\
$\beta_\pi$ ($\pi-\pi^*$) & $0.02$\,{\scriptsize(1.1)} & $0.02$\,{\scriptsize(1.1)} & $0.03$\,{\scriptsize(1.6)} \\
$\Delta\beta_\pi^{\text{High}}$ & $0.01$\,{\scriptsize(0.7)} & $0.03$\,{\scriptsize(1.5)} & $0.01$\,{\scriptsize(0.5)} \\
$\beta_x$ (ugap) & $-0.03^{***}$\,{\scriptsize(-2.6)} & $-0.02^{**}$\,{\scriptsize(-2.1)} & $-0.03^{**}$\,{\scriptsize(-2.5)} \\
\addlinespace
$n$ & 459 & 398 & 278 \\
high-conn.\ share & 0.23 & 0.21 & 0.12 \\
long-run $\beta_\pi$ (low/high) & $1.01/1.73$ & $1.42/3.08$ & $1.48/2.05$ \\
\midrule
\multicolumn{4}{l}{\textit{$\tau=0.7$}} \\
Intercept & $0.06^{*}$\,{\scriptsize(1.8)} & $0.04$\,{\scriptsize(1.2)} & $0.05$\,{\scriptsize(1.4)} \\
$\rho$ ($i_{t-1}$) & $0.97^{***}$\,{\scriptsize(125.9)} & $0.98^{***}$\,{\scriptsize(111.2)} & $0.97^{***}$\,{\scriptsize(94.9)} \\
$\beta_\pi$ ($\pi-\pi^*$) & $0.01$\,{\scriptsize(0.8)} & $0.02$\,{\scriptsize(1.3)} & $0.02$\,{\scriptsize(1.2)} \\
$\Delta\beta_\pi^{\text{High}}$ & $0.05^{**}$\,{\scriptsize(2.6)} & $0.04^{**}$\,{\scriptsize(2.3)} & $0.06^{***}$\,{\scriptsize(3.5)} \\
$\beta_x$ (ugap) & $-0.03^{***}$\,{\scriptsize(-3.0)} & $-0.02^{**}$\,{\scriptsize(-2.2)} & $-0.03^{***}$\,{\scriptsize(-2.7)} \\
\addlinespace
$n$ & 459 & 398 & 278 \\
high-conn.\ share & 0.41 & 0.41 & 0.41 \\
long-run $\beta_\pi$ (low/high) & $0.46/2.41$ & $1.05/3.36$ & $0.64/2.85$ \\
\midrule
\multicolumn{4}{l}{\textit{$\tau=0.9$}} \\
Intercept & $0.06^{*}$\,{\scriptsize(1.8)} & $0.04$\,{\scriptsize(1.1)} & $0.05$\,{\scriptsize(1.4)} \\
$\rho$ ($i_{t-1}$) & $0.98^{***}$\,{\scriptsize(121.8)} & $0.98^{***}$\,{\scriptsize(110.1)} & $0.98^{***}$\,{\scriptsize(99.8)} \\
$\beta_\pi$ ($\pi-\pi^*$) & $0.01$\,{\scriptsize(0.6)} & $0.02$\,{\scriptsize(0.7)} & $0.05^{***}$\,{\scriptsize(3.3)} \\
$\Delta\beta_\pi^{\text{High}}$ & $0.03$\,{\scriptsize(1.2)} & $0.03$\,{\scriptsize(1.3)} & $-0.03$\,{\scriptsize(-1.0)} \\
$\beta_x$ (ugap) & $-0.04^{***}$\,{\scriptsize(-3.0)} & $-0.02^{**}$\,{\scriptsize(-2.0)} & $-0.03^{**}$\,{\scriptsize(-2.5)} \\
\addlinespace
$n$ & 459 & 398 & 278 \\
high-conn.\ share & 0.59 & 0.61 & 0.42 \\
long-run $\beta_\pi$ (low/high) & $0.48/1.66$ & $0.82/2.70$ & $2.08/0.92$ \\
\bottomrule
\end{tabular}
}\par\bigskip

\textbf{Panel B. Mismatched: MICH and month-on-month CPI inflation}\par\smallskip
{\setlength{\tabcolsep}{4pt}\scriptsize
\begin{tabular}{lccc}
\toprule
 & Earliest (1987) & 1993 & 2003 \\
\midrule
\multicolumn{4}{l}{\textit{$\tau=0.5$}} \\
Intercept & $0.05$\,{\scriptsize(1.4)} & $0.02$\,{\scriptsize(0.6)} & $0.03$\,{\scriptsize(0.8)} \\
$\rho$ ($i_{t-1}$) & $0.98^{***}$\,{\scriptsize(113.9)} & $0.99^{***}$\,{\scriptsize(110.0)} & $0.99^{***}$\,{\scriptsize(99.1)} \\
$\beta_\pi$ ($\pi-\pi^*$) & $0.00$\,{\scriptsize(0.1)} & $-0.00$\,{\scriptsize(-0.1)} & $0.04^{***}$\,{\scriptsize(3.3)} \\
$\Delta\beta_\pi^{\text{High}}$ & $0.03$\,{\scriptsize(1.1)} & $0.05^{*}$\,{\scriptsize(1.8)} & $-0.05^{*}$\,{\scriptsize(-1.8)} \\
$\beta_x$ (ugap) & $-0.03^{**}$\,{\scriptsize(-2.5)} & $-0.02$\,{\scriptsize(-1.6)} & $-0.03^{**}$\,{\scriptsize(-2.0)} \\
\addlinespace
$n$ & 458 & 397 & 277 \\
high-conn.\ share & 0.46 & 0.42 & 0.55 \\
long-run $\beta_\pi$ (low/high) & $0.13/1.59$ & $-0.29/4.54$ & $3.34/-0.35$ \\
\midrule
\multicolumn{4}{l}{\textit{$\tau=0.7$}} \\
Intercept & $0.05$\,{\scriptsize(1.4)} & $0.04$\,{\scriptsize(0.9)} & $0.04$\,{\scriptsize(1.1)} \\
$\rho$ ($i_{t-1}$) & $0.98^{***}$\,{\scriptsize(113.7)} & $0.98^{***}$\,{\scriptsize(109.1)} & $0.98^{***}$\,{\scriptsize(91.9)} \\
$\beta_\pi$ ($\pi-\pi^*$) & $0.04$\,{\scriptsize(1.6)} & $0.03$\,{\scriptsize(1.1)} & $0.03$\,{\scriptsize(1.4)} \\
$\Delta\beta_\pi^{\text{High}}$ & $-0.01$\,{\scriptsize(-0.5)} & $0.01$\,{\scriptsize(0.2)} & $0.00$\,{\scriptsize(0.1)} \\
$\beta_x$ (ugap) & $-0.03^{**}$\,{\scriptsize(-2.4)} & $-0.02^{*}$\,{\scriptsize(-1.9)} & $-0.03^{**}$\,{\scriptsize(-2.2)} \\
\addlinespace
$n$ & 458 & 397 & 277 \\
high-conn.\ share & 0.47 & 0.39 & 0.30 \\
long-run $\beta_\pi$ (low/high) & $1.71/1.02$ & $1.74/2.22$ & $1.60/1.74$ \\
\midrule
\multicolumn{4}{l}{\textit{$\tau=0.9$}} \\
Intercept & $0.05$\,{\scriptsize(1.4)} & $0.04$\,{\scriptsize(0.9)} & $0.05$\,{\scriptsize(1.2)} \\
$\rho$ ($i_{t-1}$) & $0.98^{***}$\,{\scriptsize(117.2)} & $0.98^{***}$\,{\scriptsize(105.9)} & $0.98^{***}$\,{\scriptsize(93.3)} \\
$\beta_\pi$ ($\pi-\pi^*$) & $0.03$\,{\scriptsize(1.6)} & $0.03$\,{\scriptsize(1.3)} & $0.04^{*}$\,{\scriptsize(1.7)} \\
$\Delta\beta_\pi^{\text{High}}$ & $-0.00$\,{\scriptsize(-0.2)} & $0.02$\,{\scriptsize(1.0)} & $-0.00$\,{\scriptsize(-0.2)} \\
$\beta_x$ (ugap) & $-0.03^{***}$\,{\scriptsize(-2.6)} & $-0.02^{*}$\,{\scriptsize(-2.0)} & $-0.03^{**}$\,{\scriptsize(-2.2)} \\
\addlinespace
$n$ & 458 & 397 & 277 \\
high-conn.\ share & 0.40 & 0.35 & 0.36 \\
long-run $\beta_\pi$ (low/high) & $1.34/1.15$ & $1.57/2.57$ & $1.73/1.58$ \\
\bottomrule
\end{tabular}
}

\begin{tablenotes}\footnotesize
\item \textit{Notes:} Estimates of equation~\eqref{eq:regime_taylor} in which the
connectedness regime is identified from the inflation--expectations system over a 120-month
rolling window. Panel~A pairs the Michigan one-year inflation expectation with 12-month CPI
inflation; Panel~B pairs it with month-on-month CPI inflation. Columns are sample start
dates; within each panel the three blocks are the quantiles $\tau=0.5$, $0.7$ and $0.9$.
Newey--West HAC $t$-statistics in parentheses; $^{*}$, $^{**}$, $^{***}$ denote significance
at 10, 5 and 1 per cent. $\Delta\beta_\pi^{\text{High}}$ is the interaction $\delta$; the
final row of each block reports the implied long-run response $\beta_\pi/(1-\rho)$ in the
low and high regimes.
\end{tablenotes}
\end{threeparttable}
\end{table}

A limitation of the preceding analysis is that, at a monthly frequency, connectedness is
estimated with only two lags, so lagged spillovers are captured over a two-month horizon.
This restriction is necessary: within the short rolling windows the underlying vector
autoregression must remain estimable, and the implied variance-decomposition matrix must be
well-conditioned and positive semi-definite. Additional lags rapidly exhaust the available
degrees of freedom and degrade the connectedness estimates. The concern, however, is that the
price spirals and broad subcomponent co-movements that might concern policymakers could build
over a horizon longer than two months. We therefore re-estimate connectedness at a quarterly
frequency, again with two lags, which extends the horizon over which spillovers can propagate
to roughly six months. The Taylor rule continues to be estimated at a monthly frequency; a
given month is assigned to the high- or low-connectedness regime according to the state of
the quarter in which it falls.


Table~\ref{tab:sixcpiQ} reports the Taylor-rule estimates for the six CPI subcomponents when
the connectedness regime is based on connecedness measured at quarterly frequency. The evidence is notably stronger than under the monthly measure. Of the eight estimates of the interaction $\delta$ at
$\tau=0.5$ and $\tau=0.7$---four sample start dates at each quantile---six are positive and
statistically significant at the 10 per cent level or better (most at 5 per cent or better),
indicating a stronger inflation response in high-connectedness months. At $\tau=0.9$, by
contrast, only one of the four estimates is statistically significant.

\begin{table}[!htbp]
\centering
\caption{Subcomponent connectedness: Quarterly }
\label{tab:sixcpiQ}

 
 

\setlength{\tabcolsep}{4pt}\scriptsize
\begin{tabular}{lcccc}
\toprule
 & Earliest (1977) & 1983 & 1993 & 2003 \\
\midrule
\multicolumn{5}{l}{\textit{$\tau=0.5$}} \\
Intercept & $0.17^{**}$\,{\scriptsize(2.2)} & $0.12^{*}$\,{\scriptsize(1.9)} & $0.03$\,{\scriptsize(0.9)} & $0.04$\,{\scriptsize(1.2)} \\
$\rho$ ($i_{t-1}$) & $0.95^{***}$\,{\scriptsize(52.4)} & $0.96^{***}$\,{\scriptsize(61.7)} & $0.99^{***}$\,{\scriptsize(132.4)} & $0.98^{***}$\,{\scriptsize(112.7)} \\
$\beta_\pi$ ($\pi-\pi^*$) & $0.05^{**}$\,{\scriptsize(2.3)} & $0.01$\,{\scriptsize(0.4)} & $0.01$\,{\scriptsize(0.4)} & $0.01$\,{\scriptsize(0.6)} \\
$\Delta\beta_\pi^{\text{High}}$ & $0.03$\,{\scriptsize(1.5)} & $0.06^{***}$\,{\scriptsize(2.7)} & $0.06^{***}$\,{\scriptsize(3.2)} & $0.05^{***}$\,{\scriptsize(3.0)} \\
$\beta_x$ (ugap) & $-0.05^{**}$\,{\scriptsize(-2.6)} & $-0.04^{**}$\,{\scriptsize(-2.4)} & $-0.02^{**}$\,{\scriptsize(-2.3)} & $-0.03^{***}$\,{\scriptsize(-3.0)} \\
\addlinespace
$n$ (months) & 589 & 517 & 397 & 277 \\
high-conn.\ share & 0.32 & 0.28 & 0.13 & 0.19 \\
long-run $\beta_\pi$ (low/high) & $0.91/1.42$ & $0.14/1.55$ & $0.38/4.41$ & $0.39/2.85$ \\
\midrule
\multicolumn{5}{l}{\textit{$\tau=0.7$}} \\
Intercept & $0.16^{**}$\,{\scriptsize(2.1)} & $0.11^{*}$\,{\scriptsize(1.8)} & $0.02$\,{\scriptsize(0.5)} & $0.03$\,{\scriptsize(0.8)} \\
$\rho$ ($i_{t-1}$) & $0.95^{***}$\,{\scriptsize(54.2)} & $0.96^{***}$\,{\scriptsize(61.1)} & $0.99^{***}$\,{\scriptsize(127.8)} & $0.98^{***}$\,{\scriptsize(109.5)} \\
$\beta_\pi$ ($\pi-\pi^*$) & $0.02$\,{\scriptsize(1.3)} & $0.01$\,{\scriptsize(0.7)} & $0.00$\,{\scriptsize(0.1)} & $0.01$\,{\scriptsize(0.5)} \\
$\Delta\beta_\pi^{\text{High}}$ & $0.05^{**}$\,{\scriptsize(2.2)} & $0.03$\,{\scriptsize(1.4)} & $0.05^{**}$\,{\scriptsize(2.2)} & $0.04^{*}$\,{\scriptsize(1.8)} \\
$\beta_x$ (ugap) & $-0.05^{***}$\,{\scriptsize(-2.6)} & $-0.04^{**}$\,{\scriptsize(-2.4)} & $-0.02^{**}$\,{\scriptsize(-2.0)} & $-0.02^{**}$\,{\scriptsize(-2.1)} \\
\addlinespace
$n$ (months) & 589 & 517 & 397 & 277 \\
high-conn.\ share & 0.42 & 0.34 & 0.20 & 0.31 \\
long-run $\beta_\pi$ (low/high) & $0.46/1.35$ & $0.37/1.24$ & $0.14/5.14$ & $0.60/3.06$ \\
\midrule
\multicolumn{5}{l}{\textit{$\tau=0.9$}} \\
Intercept & $0.16^{**}$\,{\scriptsize(2.4)} & $0.11^{*}$\,{\scriptsize(1.7)} & $0.03$\,{\scriptsize(0.7)} & $0.03$\,{\scriptsize(1.0)} \\
$\rho$ ($i_{t-1}$) & $0.95^{***}$\,{\scriptsize(57.9)} & $0.96^{***}$\,{\scriptsize(59.3)} & $0.99^{***}$\,{\scriptsize(111.0)} & $0.98^{***}$\,{\scriptsize(105.0)} \\
$\beta_\pi$ ($\pi-\pi^*$) & $0.07^{***}$\,{\scriptsize(3.5)} & $0.02$\,{\scriptsize(0.7)} & $0.01$\,{\scriptsize(0.6)} & $0.01$\,{\scriptsize(0.5)} \\
$\Delta\beta_\pi^{\text{High}}$ & $-0.03$\,{\scriptsize(-1.4)} & $0.03$\,{\scriptsize(1.3)} & $0.03$\,{\scriptsize(1.3)} & $0.04^{*}$\,{\scriptsize(1.7)} \\
$\beta_x$ (ugap) & $-0.04^{**}$\,{\scriptsize(-2.3)} & $-0.04^{**}$\,{\scriptsize(-2.0)} & $-0.02^{*}$\,{\scriptsize(-1.7)} & $-0.02^{**}$\,{\scriptsize(-2.3)} \\
\addlinespace
$n$ (months) & 589 & 517 & 397 & 277 \\
high-conn.\ share & 0.40 & 0.44 & 0.39 & 0.36 \\
long-run $\beta_\pi$ (low/high) & $1.41/0.89$ & $0.40/1.25$ & $0.83/3.13$ & $0.60/2.86$ \\
\bottomrule
\end{tabular}

\begin{tablenotes}\footnotesize
\item \textit{Notes:} Estimates of equation~\eqref{eq:regime_taylor}. Columns are sample
start dates; each panel a quantile $\tau$. Newey--West HAC $t$-statistics in parentheses;
$^{*}$, $^{**}$, $^{***}$ denote significance at 10, 5 and 1 per cent. $\Delta\beta_\pi^{
\text{High}}$ is the interaction $\delta$; the final row reports the implied long-run
response $\beta_\pi/(1-\rho)$ in the low and high regimes.
\end{tablenotes}
\end{table}

Table~\ref{tab:expconnQ} repeats the analysis for the inflation--expectations system under
quarterly connectedness; an added advantage of this is quarterly inflation expectations are available back to 1960. Using year-on-year inflation, matched to the one-year horizon of
the Michigan survey, any evidence of a stronger inflation response is confined to
$\tau=0.5$, where the interaction $\delta$ is positive and statistically significant in the
1993 and 2003 samples. At $\tau=0.7$ and $\tau=0.9$ it is statistically insignificant, and in
one specification significantly negative. Aggregating connectedness to a quarterly frequency
therefore does not strengthen the evidence that inflation--expectations connectedness shifts
the monetary-policy reaction function.\footnote{This conclusion holds if inflation is measured as quarterly percentage change as well}.

\begin{table}[!htbp]
\centering
\caption{Expectations connectedness: Quarterly }
\label{tab:expcpiQ}

 
 

\setlength{\tabcolsep}{4pt}\scriptsize
\begin{tabular}{lcccc}
\toprule
 & Extended$^{\dagger}$ (1969) & Earliest (1987) & 1993 & 2003 \\
\midrule
\multicolumn{5}{l}{\textit{$\tau=0.5$}} \\
Intercept & $0.20^{***}$\,{\scriptsize(2.9)} & $0.05$\,{\scriptsize(1.3)} & $0.03$\,{\scriptsize(0.8)} & $0.04$\,{\scriptsize(1.1)} \\
$\rho$ ($i_{t-1}$) & $0.94^{***}$\,{\scriptsize(53.8)} & $0.98^{***}$\,{\scriptsize(109.4)} & $0.99^{***}$\,{\scriptsize(118.7)} & $0.98^{***}$\,{\scriptsize(112.6)} \\
$\beta_\pi$ ($\pi-\pi^*$) & $0.08^{***}$\,{\scriptsize(3.2)} & $0.02$\,{\scriptsize(1.3)} & $0.01$\,{\scriptsize(0.6)} & $0.01$\,{\scriptsize(0.3)} \\
$\Delta\beta_\pi^{\text{High}}$ & $-0.05^{**}$\,{\scriptsize(-2.3)} & $0.01$\,{\scriptsize(0.6)} & $0.05^{***}$\,{\scriptsize(3.0)} & $0.04^{*}$\,{\scriptsize(1.9)} \\
$\beta_x$ (ugap) & $-0.06^{***}$\,{\scriptsize(-2.8)} & $-0.03^{**}$\,{\scriptsize(-2.5)} & $-0.02^{**}$\,{\scriptsize(-2.1)} & $-0.03^{***}$\,{\scriptsize(-2.8)} \\
\addlinespace
$n$ (months) & 674 & 460 & 397 & 277 \\
high-conn.\ share & 0.16 & 0.27 & 0.31 & 0.42 \\
long-run $\beta_\pi$ (low/high) & $1.26/0.40$ & $1.04/1.66$ & $0.61/3.88$ & $0.30/2.61$ \\
\midrule
\multicolumn{5}{l}{\textit{$\tau=0.7$}} \\
Intercept & $0.16^{***}$\,{\scriptsize(2.8)} & $0.06$\,{\scriptsize(1.5)} & $0.05$\,{\scriptsize(1.3)} & $0.05$\,{\scriptsize(1.3)} \\
$\rho$ ($i_{t-1}$) & $0.95^{***}$\,{\scriptsize(67.1)} & $0.98^{***}$\,{\scriptsize(104.8)} & $0.98^{***}$\,{\scriptsize(105.2)} & $0.97^{***}$\,{\scriptsize(90.7)} \\
$\beta_\pi$ ($\pi-\pi^*$) & $0.03$\,{\scriptsize(1.0)} & $0.03$\,{\scriptsize(1.6)} & $0.04^{***}$\,{\scriptsize(2.7)} & $0.01$\,{\scriptsize(0.3)} \\
$\Delta\beta_\pi^{\text{High}}$ & $0.07^{**}$\,{\scriptsize(2.5)} & $-0.01$\,{\scriptsize(-0.2)} & $-0.05^{*}$\,{\scriptsize(-1.8)} & $0.03$\,{\scriptsize(1.6)} \\
$\beta_x$ (ugap) & $-0.05^{***}$\,{\scriptsize(-2.7)} & $-0.03^{**}$\,{\scriptsize(-2.6)} & $-0.03^{**}$\,{\scriptsize(-2.1)} & $-0.03^{**}$\,{\scriptsize(-2.2)} \\
\addlinespace
$n$ (months) & 674 & 460 & 397 & 277 \\
high-conn.\ share & 0.56 & 0.22 & 0.21 & 0.90 \\
long-run $\beta_\pi$ (low/high) & $0.55/1.79$ & $1.28/0.95$ & $2.11/-0.47$ & $0.24/1.59$ \\
\midrule
\multicolumn{5}{l}{\textit{$\tau=0.9$}} \\
Intercept & $0.18^{***}$\,{\scriptsize(2.9)} & $0.05$\,{\scriptsize(1.5)} & $0.04$\,{\scriptsize(0.9)} & $0.04$\,{\scriptsize(1.2)} \\
$\rho$ ($i_{t-1}$) & $0.95^{***}$\,{\scriptsize(64.8)} & $0.98^{***}$\,{\scriptsize(109.3)} & $0.98^{***}$\,{\scriptsize(101.3)} & $0.98^{***}$\,{\scriptsize(97.0)} \\
$\beta_\pi$ ($\pi-\pi^*$) & $0.05$\,{\scriptsize(1.3)} & $0.03$\,{\scriptsize(1.4)} & $0.02$\,{\scriptsize(0.9)} & $0.01$\,{\scriptsize(0.6)} \\
$\Delta\beta_\pi^{\text{High}}$ & $0.01$\,{\scriptsize(0.2)} & $-0.01$\,{\scriptsize(-0.2)} & $0.01$\,{\scriptsize(0.4)} & $0.03$\,{\scriptsize(1.0)} \\
$\beta_x$ (ugap) & $-0.05^{***}$\,{\scriptsize(-2.9)} & $-0.03^{**}$\,{\scriptsize(-2.5)} & $-0.02^{*}$\,{\scriptsize(-1.9)} & $-0.03^{**}$\,{\scriptsize(-2.6)} \\
\addlinespace
$n$ (months) & 674 & 460 & 397 & 277 \\
high-conn.\ share & 0.64 & 0.30 & 0.27 & 0.24 \\
long-run $\beta_\pi$ (low/high) & $0.92/1.08$ & $1.39/1.13$ & $1.53/2.25$ & $0.84/2.39$ \\
\bottomrule
\end{tabular}

\begin{tablenotes}\footnotesize
\item \textit{Notes:} Estimates of equation~\eqref{eq:regime_taylor}. Columns are sample
start dates; each panel a quantile $\tau$. Newey--West HAC $t$-statistics in parentheses;
$^{*}$, $^{**}$, $^{***}$ denote significance at 10, 5 and 1 per cent. $\Delta\beta_\pi^{
\text{High}}$ is the interaction $\delta$; the final row reports the implied long-run
response $\beta_\pi/(1-\rho)$ in the low and high regimes. $^{\dagger}$ The Extended column uses the quarterly Michigan survey (available from 1960); the remaining columns use the monthly survey (from 1978) aggregated to quarters.
\end{tablenotes}
\end{table}
 
% ============================================================
\appendix
\section{Connectedness matrices}
% ============================================================
Tables~\ref{tab:conncfull}--\ref{tab:connlcore} report the full pairwise connectedness structure
among the six CPI subcomponents at $\tau=0.9$ ($n_{\text{lag}}=2$), separately for the
contemporaneous and lagged blocks, for the full sample and the two subperiods (1967--1982 and
1983--2026). In each matrix, cell $(i,j)$ is the share of row $i$'s variation attributed to column
$j$; FROM is the row sum (received), TO the column sum (transmitted), and NET $=$ TO$-$FROM. The
contemporaneous block is close to symmetric (small NET), whereas the lagged block carries the
directional structure --- shelter and core services are the net lagged transmitters. Gasoline is
weakly connected to the other categories: its variation is dominated by energy-market (crude-oil)
movements that lie outside the CPI system, so it sits at the periphery of the network.

\begin{table}[!ht]
\centering
\caption{Contemporaneous connectedness among CPI subcomponents (full sample, 1967--2026, $\tau=0.9$)}
\label{tab:conncfull}
\setlength{\tabcolsep}{4pt}\footnotesize
\begin{threeparttable}
\begin{tabular}{l c c c c c c c}
\toprule
 & Food & Gas. & HH en. & Core gds & Shelter & Core svc. & FROM \\
\midrule
Food & --- & 0.2 & 1.6 & 4.6 & 5.4 & 1.2 & 12.9 \\
Gasoline & 0.4 & --- & 4.9 & 0.1 & 0.1 & 0.1 & 5.6 \\
HH energy & 1.6 & 4.4 & --- & 1.8 & 1.0 & 1.2 & 9.9 \\
Core goods & 4.6 & 0.0 & 1.8 & --- & 6.4 & 2.5 & 15.3 \\
Shelter & 4.9 & 0.0 & 0.5 & 6.3 & --- & 5.8 & 17.5 \\
Core svc.\ ex-sh. & 0.8 & 0.1 & 1.2 & 2.3 & 5.7 & --- & 10.1 \\
\midrule
TO & 12.2 & 4.7 & 10.0 & 15.2 & 18.5 & 10.8 &  \\
NET & -0.7 & -0.9 & +0.0 & -0.2 & +1.0 & +0.7 &  \\
\bottomrule
\end{tabular}
\begin{tablenotes}[flushleft]\footnotesize
\item \textit{Notes:} Pseudo-quantile $R^2$ connectedness (\%) among the six CPI subcomponents at $\tau=0.9$, $n_{\text{lag}}=2$. Cell $(i,j)$ is the share of row $i$'s variation attributed to column $j$. FROM $=$ row sum (received); TO $=$ column sum (transmitted); NET $=$ TO$-$FROM ($>0$ net transmitter). Diagonal (own) terms are excluded from FROM/TO.
\end{tablenotes}
\end{threeparttable}
\end{table}


\begin{table}[!ht]
\centering
\caption{Lagged connectedness among CPI subcomponents (full sample, 1967--2026, $\tau=0.9$)}
\label{tab:connlfull}
\setlength{\tabcolsep}{4pt}\footnotesize
\begin{threeparttable}
\begin{tabular}{l c c c c c c c}
\toprule
 & Food & Gas. & HH en. & Core gds & Shelter & Core svc. & FROM \\
\midrule
Food & 23.5 & 0.3 & 5.5 & 6.6 & 13.1 & 4.6 & 30.1 \\
Gasoline & 1.4 & 11.3 & 1.4 & 0.3 & 0.3 & 0.3 & 3.6 \\
HH energy & 2.8 & 2.8 & 10.2 & 2.6 & 1.9 & 6.2 & 16.3 \\
Core goods & 6.8 & 0.4 & 3.1 & 29.3 & 5.6 & 9.9 & 25.8 \\
Shelter & 12.0 & 0.1 & 1.3 & 6.5 & 26.6 & 12.8 & 32.7 \\
Core svc.\ ex-sh. & 3.5 & 0.9 & 3.6 & 7.1 & 16.2 & 15.4 & 31.4 \\
\midrule
TO & 26.5 & 4.5 & 14.9 & 23.1 & 37.1 & 33.8 &  \\
NET & -3.6 & +0.9 & -1.4 & -2.8 & +4.4 & +2.4 &  \\
\bottomrule
\end{tabular}
\begin{tablenotes}[flushleft]\footnotesize
\item \textit{Notes:} Pseudo-quantile $R^2$ connectedness (\%) among the six CPI subcomponents at $\tau=0.9$, $n_{\text{lag}}=2$. Cell $(i,j)$ is the share of row $i$'s variation attributed to column $j$. FROM $=$ row sum (received); TO $=$ column sum (transmitted); NET $=$ TO$-$FROM ($>0$ net transmitter). Diagonal (own) terms are excluded from FROM/TO (and shown for the lagged block as the own-lag/persistence term).
\end{tablenotes}
\end{threeparttable}
\end{table}


\begin{table}[!ht]
\centering
\caption{Contemporaneous connectedness among CPI subcomponents (1967--1982, $\tau=0.9$)}
\label{tab:conncgi}
\setlength{\tabcolsep}{4pt}\footnotesize
\begin{threeparttable}
\begin{tabular}{l c c c c c c c}
\toprule
 & Food & Gas. & HH en. & Core gds & Shelter & Core svc. & FROM \\
\midrule
Food & --- & 4.9 & 1.9 & 8.5 & 2.9 & 1.3 & 19.6 \\
Gasoline & 3.6 & --- & 6.9 & 4.0 & 4.7 & 7.9 & 27.0 \\
HH energy & 1.5 & 18.9 & --- & 3.1 & 1.6 & 2.8 & 28.1 \\
Core goods & 4.3 & 11.5 & 2.3 & --- & 4.6 & 4.8 & 27.5 \\
Shelter & 2.2 & 16.2 & 2.0 & 4.8 & --- & 10.5 & 35.8 \\
Core svc.\ ex-sh. & 2.9 & 16.6 & 1.9 & 4.1 & 5.4 & --- & 31.0 \\
\midrule
TO & 14.6 & 68.2 & 15.0 & 24.6 & 19.3 & 27.3 &  \\
NET & -5.0 & +41.2 & -13.1 & -2.9 & -16.5 & -3.7 &  \\
\bottomrule
\end{tabular}
\begin{tablenotes}[flushleft]\footnotesize
\item \textit{Notes:} Pseudo-quantile $R^2$ connectedness (\%) among the six CPI subcomponents at $\tau=0.9$, $n_{\text{lag}}=2$. Cell $(i,j)$ is the share of row $i$'s variation attributed to column $j$. FROM $=$ row sum (received); TO $=$ column sum (transmitted); NET $=$ TO$-$FROM ($>0$ net transmitter). Diagonal (own) terms are excluded from FROM/TO.
\end{tablenotes}
\end{threeparttable}
\end{table}


\begin{table}[!ht]
\centering
\caption{Lagged connectedness among CPI subcomponents (1967--1982, $\tau=0.9$)}
\label{tab:connlgi}
\setlength{\tabcolsep}{4pt}\footnotesize
\begin{threeparttable}
\begin{tabular}{l c c c c c c c}
\toprule
 & Food & Gas. & HH en. & Core gds & Shelter & Core svc. & FROM \\
\midrule
Food & 31.5 & 2.1 & 23.2 & 9.4 & 3.8 & 6.0 & 44.4 \\
Gasoline & 22.3 & 24.4 & 10.9 & 4.3 & 8.4 & 2.8 & 48.6 \\
HH energy & 15.8 & 8.1 & 15.9 & 19.7 & 4.8 & 7.3 & 55.8 \\
Core goods & 6.9 & 4.8 & 17.4 & 34.7 & 5.0 & 3.4 & 37.6 \\
Shelter & 15.5 & 6.3 & 2.8 & 20.4 & 16.5 & 2.5 & 47.5 \\
Core svc.\ ex-sh. & 16.2 & 9.5 & 3.6 & 18.6 & 5.3 & 15.5 & 53.3 \\
\midrule
TO & 76.7 & 30.9 & 57.9 & 72.4 & 27.4 & 21.9 &  \\
NET & +32.3 & -17.7 & +2.1 & +34.9 & -20.1 & -31.4 &  \\
\bottomrule
\end{tabular}
\begin{tablenotes}[flushleft]\footnotesize
\item \textit{Notes:} Pseudo-quantile $R^2$ connectedness (\%) among the six CPI subcomponents at $\tau=0.9$, $n_{\text{lag}}=2$. Cell $(i,j)$ is the share of row $i$'s variation attributed to column $j$. FROM $=$ row sum (received); TO $=$ column sum (transmitted); NET $=$ TO$-$FROM ($>0$ net transmitter). Diagonal (own) terms are excluded from FROM/TO (and shown for the lagged block as the own-lag/persistence term).
\end{tablenotes}
\end{threeparttable}
\end{table}


\begin{table}[!ht]
\centering
\caption{Contemporaneous connectedness among CPI subcomponents (1983--2026, $\tau=0.9$)}
\label{tab:connccore}
\setlength{\tabcolsep}{4pt}\footnotesize
\begin{threeparttable}
\begin{tabular}{l c c c c c c c}
\toprule
 & Food & Gas. & HH en. & Core gds & Shelter & Core svc. & FROM \\
\midrule
Food & --- & 0.2 & 6.3 & 0.3 & 2.2 & 1.0 & 9.9 \\
Gasoline & 0.4 & --- & 4.0 & 0.1 & 3.5 & 0.3 & 8.2 \\
HH energy & 6.5 & 3.8 & --- & 1.1 & 0.2 & 0.4 & 12.0 \\
Core goods & 0.3 & 0.1 & 1.0 & --- & 0.4 & 4.1 & 5.8 \\
Shelter & 2.2 & 2.9 & 0.2 & 0.4 & --- & 2.8 & 8.5 \\
Core svc.\ ex-sh. & 1.0 & 0.2 & 0.4 & 4.2 & 2.8 & --- & 8.6 \\
\midrule
TO & 10.4 & 7.2 & 11.8 & 6.0 & 9.1 & 8.4 &  \\
NET & +0.5 & -1.0 & -0.2 & +0.2 & +0.6 & -0.1 &  \\
\bottomrule
\end{tabular}
\begin{tablenotes}[flushleft]\footnotesize
\item \textit{Notes:} Pseudo-quantile $R^2$ connectedness (\%) among the six CPI subcomponents at $\tau=0.9$, $n_{\text{lag}}=2$. Cell $(i,j)$ is the share of row $i$'s variation attributed to column $j$. FROM $=$ row sum (received); TO $=$ column sum (transmitted); NET $=$ TO$-$FROM ($>0$ net transmitter). Diagonal (own) terms are excluded from FROM/TO.
\end{tablenotes}
\end{threeparttable}
\end{table}


\begin{table}[!ht]
\centering
\caption{Lagged connectedness among CPI subcomponents (1983--2026, $\tau=0.9$)}
\label{tab:connlcore}
\setlength{\tabcolsep}{4pt}\footnotesize
\begin{threeparttable}
\begin{tabular}{l c c c c c c c}
\toprule
 & Food & Gas. & HH en. & Core gds & Shelter & Core svc. & FROM \\
\midrule
Food & 14.6 & 1.6 & 9.2 & 2.6 & 4.8 & 0.7 & 18.9 \\
Gasoline & 3.6 & 10.3 & 3.0 & 0.0 & 0.8 & 0.6 & 8.0 \\
HH energy & 4.6 & 3.8 & 10.8 & 0.8 & 0.8 & 1.2 & 11.2 \\
Core goods & 2.6 & 1.6 & 1.4 & 30.7 & 0.7 & 3.0 & 9.3 \\
Shelter & 7.5 & 0.5 & 0.5 & 0.5 & 19.9 & 8.3 & 17.4 \\
Core svc.\ ex-sh. & 3.1 & 8.1 & 3.1 & 2.1 & 8.8 & 5.7 & 25.1 \\
\midrule
TO & 21.4 & 15.6 & 17.3 & 6.0 & 15.8 & 13.9 &  \\
NET & +2.4 & +7.6 & +6.1 & -3.3 & -1.6 & -11.3 &  \\
\bottomrule
\end{tabular}
\begin{tablenotes}[flushleft]\footnotesize
\item \textit{Notes:} Pseudo-quantile $R^2$ connectedness (\%) among the six CPI subcomponents at $\tau=0.9$, $n_{\text{lag}}=2$. Cell $(i,j)$ is the share of row $i$'s variation attributed to column $j$. FROM $=$ row sum (received); TO $=$ column sum (transmitted); NET $=$ TO$-$FROM ($>0$ net transmitter). Diagonal (own) terms are excluded from FROM/TO (and shown for the lagged block as the own-lag/persistence term).
\end{tablenotes}
\end{threeparttable}
\end{table}



% ============================================================
\section{Goods-to-services sequencing: a lag-window illustration}
% ============================================================
A frequently cited feature of the 2021--22 surge is that goods inflation led services inflation.
Our framework recovers this, but only once the measurement window is long enough --- which is itself
instructive about connectedness measurement. Table~\ref{tab:h5} reports, on the post-1983 sample at
$\tau=0.9$, the net lagged connectedness from core goods to core services. Under the monthly two-lag
dependence (Genizi) metric the sign is slightly negative ($-1.2$): a two-month window is too short
to span a multi-month goods-then-services sequence. The LASSO predictive metric, which admits longer
lags, recovers the ordering ($+1.3$ at two lags, $+5.9$ at six, where core goods becomes a net
transmitter). Table~\ref{tab:h5q} confirms it at quarterly frequency, where a two-lag block spans
six months: the lagged goods$\to$services term then exceeds the reverse at $\tau=0.5$ ($9.9$ vs
$4.4$) and $\tau=0.7$ ($11.2$ vs $7.1$). The point is methodological --- lag-window length
materially affects what connectedness detects, which motivates the quarterly analysis used elsewhere
--- rather than a new fact about the episode.

\begin{table}[!ht]
\centering
\caption{Goods-to-services sequencing under alternative metrics and lag lengths}
\label{tab:h5}
\setlength{\tabcolsep}{6pt}\footnotesize
\begin{threeparttable}
\begin{tabular}{l c c}
\toprule
Metric & Goods$\to$services (lagged, net) & Core goods NET \\
\midrule
Genizi dependence ($n_{\text{lag}}{=}2$) & -1.2 & -3.4 \\
LASSO predictive ($n_{\text{lag}}{=}2$) & +1.3 & +3.4 \\
LASSO predictive ($n_{\text{lag}}{=}6$) & +5.9 & +12.9 \\
\bottomrule
\end{tabular}
\begin{tablenotes}[flushleft]\footnotesize
\item \textit{Notes:} ``Goods$\to$services (lagged, net)'' is lagged connectedness from core goods to core services (shelter $+$ core services ex-shelter) minus the reverse, at $\tau=0.9$ on the ex-Great-Inflation sample (1983--2026); $>0$ means goods lead. The two-lag Genizi dependence metric cannot see a multi-month sequence and tilts negative; the LASSO predictive metric confirms the sequencing, increasingly at longer lags. Predictive ordering, not identified causation.
\end{tablenotes}
\end{threeparttable}
\end{table}


\begin{table}[!ht]
\centering
\caption{Goods-to-services: contemporaneous and lagged connectedness (quarterly)}
\label{tab:h5q}
\setlength{\tabcolsep}{6pt}\footnotesize
\begin{threeparttable}
\begin{tabular}{c c c c c}
\toprule
 & \multicolumn{2}{c}{Goods $\to$ Services} & \multicolumn{2}{c}{Services $\to$ Goods} \\
\cmidrule(lr){2-3}\cmidrule(lr){4-5}
$\tau$ & Contemp. & Lagged & Contemp. & Lagged \\
\midrule
0.5 & 3.3 & 9.9 & 3.4 & 4.4 \\
0.7 & 8.5 & 11.2 & 8.7 & 7.1 \\
0.9 & 13.1 & 18.4 & 15.2 & 18.3 \\
\bottomrule
\end{tabular}
\begin{tablenotes}[flushleft]\footnotesize
\item \textit{Notes:} Pairwise pseudo-quantile $R^2$ connectedness (\%) between core goods and core services (shelter $+$ core services ex-shelter) on the ex-Great-Inflation quarterly sample (1983Q1--2026Q2, 174 quarters), $n_{\text{lag}}=2$ (a six-month lag horizon). ``Goods $\to$ Services'' is the share of services' variation explained by goods, split into contemporaneous and lagged; ``Services $\to$ Goods'' is the reverse. The \emph{lagged} Goods $\to$ Services term is the ``goods lead'' propagation channel; it exceeds the lagged Services $\to$ Goods term at every quantile (by 4.0 at $\tau=0.7$), consistent with goods leading services. Six months is the horizon the monthly two-lag metric could not span. Dependence, not identified causation.
\end{tablenotes}
\end{threeparttable}
\end{table}



\printbibliography

\end{document}


